% 3_1_data.tex — Section 3.1 Data (revised per supervisor comments)

\section{Data}
\subsection{Intraday Price Data}

The primary dataset comprises high-frequency intraday price observations for the S\&P 500 index (SPX), sourced from the Bloomberg Terminal, accessed via the Quinn School of Business, University College Dublin. Data are sampled at five-minute intervals across the standard U.S. equity trading session, spanning 09:30 to 16:00 Eastern Time (ET), yielding 78 intraday return observations per trading day. The sample is restricted to complete trading sessions. Days on which U.S. equity markets were closed, including scheduled public holidays such as Christmas Day (25 December 2025) and New Year's Day (1 January 2026), and days subject to early market closure, such as Christmas Eve (24 December 2025) and the post-Thanksgiving session (28 November 2025), are excluded from the analysis. This filter ensures that each retained observation contributes a full complement of intraday returns, thereby producing realised variance estimates that are directly comparable across trading days and free from the downward bias that would otherwise arise from truncated sessions.

To assess the generalisability of the models across asset classes, supplementary five-minute spot price data are obtained from the Bloomberg Terminal for two additional instruments: the British pound against the U.S. dollar (GBP/USD) and gold denominated in U.S. dollars per troy ounce (XAU/USD). The GBP/USD foreign exchange market operates continuously from 22:00 Sunday to 22:00 Friday (UTC), though the analysis restricts attention to the London--New York overlap session (13:00--17:00 UTC) during which liquidity and price discovery are most concentrated. The XAU/USD spot gold market trades nearly around the clock on weekdays; for this study, the active session is defined as the COMEX trading hours of 08:20 to 13:30 ET. Because these supplementary assets trade over different session lengths relative to the SPX, the number of five-minute intervals per day differs across instruments, and daily realised variance will mechanically reflect the respective session durations, as follows from the summation of squared intraday returns described in Equation~(2) below.

In addition, a secondary S\&P 500 dataset spanning May 2023 to October 2023, sourced from Pineify, is incorporated to evaluate model performance under a markedly different market regime. This six-month window is characterised by relatively subdued realised variance and a notable absence of the acute geopolitical tensions and trade policy uncertainty that define the primary sample period. The period was selected on the basis that it captures contemporary market microstructure and institutional characteristics whilst exhibiting markedly lower realised and implied volatility relative to more recent periods. Notably, the CBOE Volatility Index (VIX) remained below 20 throughout this window, indicative of stable investor sentiment, while also not representing an unusually settled period. The inclusion of this comparative dataset facilitates an assessment of the extent to which model accuracy and predictive stability are sensitive to prevailing macroeconomic conditions and geopolitical risk, providing a more robust and generalisable evaluation of model performance across distinct market regimes.

% TODO: Insert figure — VIX time series (last five years) with vertical lines marking sample periods.
% \begin{figure}[htbp]
% \centering
% \includegraphics[width=\textwidth]{vix_5yr.pdf}
% \caption{CBOE Volatility Index (VIX), January 2021 to March 2026. Shaded regions indicate the primary sample period (July 2025--March 2026) and the historical comparison period (May--October 2023).}
% \label{fig:vix_5yr}
% \end{figure}

A five-minute sampling frequency is adopted in accordance with the established convention in the high-frequency finance literature. \cite{LiuPattonSheppard2015} demonstrate that this interval strikes an optimal balance between minimising microstructure noise and maximising the precision of realised variance estimates, a finding corroborated across equity, foreign exchange, and commodity markets in subsequent studies including \cite{PlihalLyocsa2021} and \cite{KambouroudisMcMillanTsakou2021}.

\subsection{Data Cleaning}

A multi-step data cleaning pipeline is applied to the raw intraday data to ensure consistency and reliability prior to realised variance construction.

All timestamps in the raw data are recorded in Irish Standard Time (IST) or Irish Summer Time (GMT+1), reflecting the timezone of the Bloomberg terminal used for extraction. Under normal conditions, when both the United States and Ireland observe their respective daylight saving time (DST) adjustments, the U.S. trading session of 09:30 to 16:00 ET corresponds to 14:30 to 21:00 Irish time. However, a complication arises from the fact that the U.S. and Ireland do not transition between standard and summer time on the same dates. During the brief windows each spring and autumn when only one jurisdiction has changed clocks, the session maps to 13:30 to 20:00 Irish time. The data cleaning pipeline identifies these DST mismatch windows programmatically and applies the appropriate session boundaries on a per-day basis to ensure that only observations recorded during active U.S. trading hours are retained.

Each observation is validated against the applicable session boundaries for its trading day, accounting for the DST adjustment described above, and records falling outside the active window are removed. Duplicate timestamps and observations with missing close prices are also dropped to prevent undefined returns from entering the realised variance calculation. U.S. market holidays and early-close sessions are then excluded. Within the sample period, six full market holidays are removed: Labor Day (1 September 2025), Thanksgiving Day (27 November 2025), Christmas Day (25 December 2025), New Year's Day (1 January 2026), Martin Luther King Jr. Day (19 January 2026), and Presidents' Day (16 February 2026). Two additional dates are excluded due to shortened trading sessions: the day after Thanksgiving (28 November 2025) and Christmas Eve (24 December 2025). Early-close days are removed entirely rather than retained with truncated sessions, as partial trading days would yield a reduced number of intraday returns and produce realised variance estimates that are not directly comparable to those from full sessions.

After cleaning, the dataset contains 164 complete trading days spanning the period from 22 July 2025 to 20 March 2026. A per-day bar count verification confirms that each retained trading day contains the expected 79 five-minute observations, with the exception of a small number of days in early March 2026 that record 78 bars due to a single missing observation. These days are retained as the impact on realised variance estimation from one missing bar is negligible.

\subsection{Macroeconomic Announcement Data}

FOMC interest rate decisions are sourced directly from the Board of Governors of the Federal Reserve System. Rather than using a survey-based consensus forecast, the expected rate outcome is derived from 30-Day Federal Funds futures prices traded on the Chicago Mercantile Exchange (CME). CME 30-Day Federal Funds futures provide a reliable measure of expected policy rates because their prices embed the aggregated expectations of market participants regarding decisions by the Federal Reserve at forthcoming FOMC meetings. This approach is consistent with \cite{BauerSwanson2023}, who demonstrate that Federal Funds futures remain reliable and efficient measures of market rate expectations around scheduled FOMC announcements.

Consensus forecast and actual release values for key U.S. jobs and inflation data that influence FOMC decisions on federal rates are sourced from the Investing.com economic calendar, a survey-based consensus forecast of professional economists and analysts at major financial institutions. Records are compiled from January 2023 onwards to construct a sufficiently large historical distribution of release outcomes, from which standard deviations are calculated for each indicator and used to implement a three-tier surprise classification system, as further described in Section 3.3. In addition to the primary Consumer Price Index (CPI) and Non-Farm Payrolls (NFP) releases, a broader set of indicators is incorporated to mitigate the constraints imposed by the limited sample size: the Producer Price Index (PPI), the Personal Consumption Expenditure (PCE) price index, the ADP National Employment Report, the Job Openings and Labour Turnover Survey (JOLTS), and Initial Jobless Claims. Each release is assigned to the three-tier surprise classification framework, and the resulting category dummies are aggregated across constituent releases to construct composite surprise indicators for inclusion in the augmented HAR-RV specifications \citep{Gilbert2017}.

For each macroeconomic indicator, the pre-computed surprise category --- recorded as an integer value of 0 (in line with expectations), 1 (moderate surprise), or 2 (large surprise) --- is used to construct per-event, per-category binary dummy variables that take the value one on the corresponding release date and zero otherwise. FOMC announcements receive separate treatment: the expected dummy (FOMC\_exp) is defined as the FOMC category-0 indicator, while the unexpected dummy (FOMC\_unexp) takes the maximum of the category-1 and category-2 indicators, thereby capturing any deviation from the futures-implied rate path. In the composite specification (Model~3a), all non-FOMC event dummies within each surprise tier are collapsed into aggregate MACRO\_0, MACRO\_1, and MACRO\_2 indicators by taking the daily maximum across constituent releases, so that the dummy equals one whenever at least one contributing indicator falls into that tier on a given trading day. A restricted variant (Model~3b) limits this aggregation to CPI and NFP releases only, excluding the secondary indicators. Dummies with zero occurrences within the in-sample window are excluded from estimation to avoid perfect multicollinearity.

\subsection{Tariff Announcement Data}

The sample period corresponds to the current episode of heightened U.S. trade policy activity, during which successive tariff announcements have generated significant market volatility. To capture these effects, a tariff event dataset is manually constructed from the Brookings Institution tariff tracker \citep{Brookings2025}, cross-referenced with contemporaneous reporting from the Wall Street Journal and Reuters. The resulting event log records 36 distinct tariff-related announcements between July 2025 and March 2026, including executive orders, retaliatory measures, legal rulings, and de-escalation agreements.

Existing approaches to quantifying trade policy uncertainty rely on text-based indices constructed from newspaper coverage \citep{BakerBloomDavis2016} or on aggregate geopolitical risk measures \citep{CaldaraIacoviello2022}. While these indices capture the general level of trade policy uncertainty over time, they do not isolate the impact of individual tariff announcements on specific trading days, nor do they distinguish between events of differing severity. \citet{HandleyLimao2017} demonstrate that trade policy uncertainty operates through discrete policy actions that alter the expected tariff path, suggesting that an event-level approach is more appropriate for capturing the immediate volatility effects of individual announcements. \citet{FajgelbaumKhandelwal2022} further document that the market impact of tariff actions varies substantially with the scope of goods covered, the number of trading partners affected, and the sectoral composition of targeted imports.

To operationalise this event-level approach, each tariff announcement is assigned a Weighted Severity Score (WSS) based on three dimensions, each scored on a scale from 0 (negligible) to 3 (severe):

\begin{itemize}
    \item \textbf{Magnitude (0--3):} The tariff rate imposed. A score of 0 indicates no meaningful rate change; 1 corresponds to rates in the 10--24\% range; 2 to rates of 25--49\%; and 3 to rates of 50\% or above.
    \item \textbf{Breadth (0--3):} The geographic scope of the action. A score of 0 denotes a single country; 1 a bilateral or small regional action; 2 a major trading bloc (e.g.\ the European Union or G20 members); and 3 a global or economy-wide measure affecting all trading partners.
    \item \textbf{Sector Impact (0--3):} The relevance of the targeted goods to the S\&P 500 index. A score of 0 indicates a negligible sector weight; 1 corresponds to a narrow sector representing less than 5\% of the index; 2 to a moderately weighted sector (5--15\%); and 3 to broad, economy-wide coverage or a sector exceeding 15\% of the index.
\end{itemize}

The WSS is computed as the unweighted sum of the three component scores, yielding a composite index ranging from 0 to 9. A binary tariff dummy (TARIFF) is then set equal to one on trading days when at least one tariff event with a WSS of five or greater occurs, and zero otherwise. The threshold of five is selected to isolate announcements of material market significance --- events that combine, at a minimum, a moderate tariff rate with either broad geographic scope or substantial sectoral relevance --- while excluding minor or procedural actions unlikely to affect aggregate equity volatility.

Of the 36 events recorded, 15 meet the WSS $\geq$ 5 threshold and are classified as material tariff events. The remaining events, which include narrow bilateral adjustments, procedural extensions, and low-magnitude measures, are excluded from the dummy. Each event is further classified as escalatory or de-escalatory based on whether the announcement represented an increase or a reduction in trade barriers, with 9 of the 36 events identified as de-escalation measures. On days when multiple tariff events occur, the maximum WSS across all events on that day determines the dummy value.

\subsection{Sample Period and Limitations}

The primary SPX sample spans 164 trading days from 22 July 2025 to 20 March 2026. The effective analysis window for most augmented model specifications begins on 19 August 2025, after the 22-day rolling window required for the monthly RV component has been populated, and extends to the end of the sample period. The historical SPX comparison dataset covers 1 May 2023 to 31 October 2023, yielding approximately 128 trading days after the removal of holidays and partial sessions. The GBP/USD and XAU/USD supplementary datasets span the same primary period as the SPX. The difference in timeframes between the primary and historical SPX samples provides an opportunity to evaluate model stability under contrasting volatility regimes: the primary period is characterised by elevated policy uncertainty and frequent tariff-related shocks, whereas the historical window reflects a comparatively stable macroeconomic environment.

For the SPX, the in-sample period runs from 22 July 2025 to 29 January 2026 and the out-of-sample period from 30 January 2026 to 20 March 2026, providing approximately 35 OOS trading days. For the historical dataset, an 80:20 in-sample to out-of-sample split is applied dynamically based on the number of complete trading days retained after cleaning. The resulting analysis window of approximately seven months for each dataset captures short-term volatility dynamics. The weekly and monthly RV components should therefore be interpreted as reflecting short-to-medium-horizon persistence rather than long-term structural behaviour, and the findings should be understood in the context of a relatively condensed sample that includes, in the primary period, an atypical episode of trade policy uncertainty.
